\documentclass[conference]{IEEEtran}
\IEEEoverridecommandlockouts
% The preceding line is only needed to identify funding in the first footnote. If that is unneeded, please comment it out.

\usepackage{cite}
\usepackage{amsmath,amssymb,amsfonts}
\usepackage{algorithmic}
\usepackage{graphicx}
\usepackage{textcomp}
\usepackage{xcolor}
\usepackage{hyperref}
\usepackage{booktabs}
\usepackage{multirow}
\usepackage{array}
\usepackage{float}

\def\BibTeX{{\rm B\kern-.05em{\sc i\kern-.025em b}\kern-.08em
    T\kern-.1667em\lower.7ex\hbox{E}\kern-.125emX}}

\begin{document}

\title{ORBITA-ATSAD: A Hybrid Database and Deep Learning Framework for Spacecraft Telemetry Anomaly Detection Using Temporal and Spatial Features}

\author{
\begin{tabular}{c@{\hspace{1.5cm}}c}
\textbf{Anik Das} & \textbf{Kshitij Chaudhari} \\
Student & Student \\
Department of Computational Intelligence & Department of Computational Intelligence \\
School of Computing & School of Computing \\
SRM Institute of Science and Technology, Kattankulathur & SRM Institute of Science and Technology, Kattankulathur \\
ad2082@srmist.edu.in & kc8488@srmist.edu.in
\end{tabular} \\[0.6cm]
\begin{tabular}{c}
\textbf{Kavitha Duraipandian} \\
Assistant Professor \\
Department of Computational Intelligence \\
School of Computing \\
SRM Institute of Science and Technology, Kattankulathur \\
kavithad7@srmist.edu.in
\end{tabular}
}


\maketitle

\begin{abstract}
This paper presents ORBITA-ATSAD, a hybrid database and deep learning framework for spacecraft telemetry anomaly detection using the ATSADBench benchmark methodology. Extending beyond simple anomaly detection, the platform serves as a comprehensive cataloging and tracking registry capable of computing orbital propagations, assessing collision risks, and managing launch vehicle data. PostgreSQL with TimescaleDB and PostGIS extensions stores the time-series telemetry data and enables 3D spatial queries for multiple satellites. The unsupervised LSTM autoencoder model, implemented in PyTorch, is trained on multivariate telemetry streams utilizing dynamic thresholding based on reconstruction error to identify patterns indicative of spacecraft failure. Integrating space weather data from NASA DONKI correlates external factors with telemetry anomalies, achieving training and validation accuracies of 89.6\% and 87.2\% respectively. The model is evaluated using ATSADBench metrics, achieving an Alarm Accuracy of 0.92, Latency of 0.87, and Contiguity of 0.89. The platform also features a dynamic benchmark registry system built with FastAPI and React, shielded by JWT authentication, and fully containerized via Docker. This multi-modal approach offers a scalable and operationally ready mechanism for spacecraft anomaly resolution.
\end{abstract}

\begin{IEEEkeywords}
Spacecraft Anomaly Detection, ATSADBench, Hybrid Database Framework, LSTM Autoencoder, TimescaleDB, PostGIS, Space Weather Correlation, React, FastAPI, Docker
\end{IEEEkeywords}

\section{INTRODUCTION}

The new technology of spacecraft telemetry monitoring has been increasing exponentially over the past few years, and today one can collect massive amounts of data from hundreds of satellites in real-time. These telemetry streams contain critical information about spacecraft health including temperature, power, attitude, and subsystem status. The analysis of this data has come to be a big issue in most space operations such as commercial satellite fleets, scientific missions, and space stations. Development of subtle anomalies that could indicate impending failures has also led to the difficulty in distinguishing between normal variations and genuine problems, and it is time to consider viable and consistent systems of detection.

The current spacecraft anomaly detectors have the tendency to lean towards either simple threshold-based alerts or post-mission forensic analysis. Even though significant strides have been achieved in the area of time-series anomaly detection using deep learning techniques such as LSTM and Transformers, applying these to operational spacecraft telemetry remains challenging due to the complexity of multivariate dependencies, varying operational modes, and external factors like space weather. In addition, the existing systems are not disposed to collaborating with the benchmark frameworks like ATSADBench which provides standardized evaluation metrics for fair comparison across different approaches.

The proposed study will fill in these gaps by coming up with a hybrid database and deep learning framework that will be capable of detecting anomalies in spacecraft telemetry streams using state-of-the-art LSTM networks evaluated against the ATSADBench benchmark. The current models that are used in the proposed architecture are: TimescaleDB for time-series telemetry storage, LSTM networks for anomaly detection, and NASA DONKI API for space weather correlation. The combination of the temporal characteristics and external environmental factors enables the system to come up with a holistic and feasible solution in detecting the spacecraft anomalies that could prevent mission failures.

This system additionally comprises a comprehensive web interface and robust backend architecture coded in React, TypeScript, FastAPI, and CesiumJS in order to make the process of visualization, cataloging, and alerting easy to use. Unlike standard detectors, the framework is integrated with an expansive tracking registry that handles astronomical orbit propagation (using SGP4/Skyfield), records real-time conjunctions, tracks orbital congestion, and logs space launch data. Securely managed through JWT authentication and containerized via Docker and Docker Compose, it provides the users with the opportunity to view satellites in 3D space via PostGIS spatial queries, register their own ATSADBench datasets dynamically as a service, and receive anomaly alerts gracefully. This type of an asynchronous hybrid solution will be developed with the intention of offering a stable, scalable, and operationally rigorous solution to detect spacecraft anomalies to improve the reliability and safety of space missions.

\subsection{Objective of the Study}

The objective of the study is to possess a state of art and hybrid database and deep learning model that with the assistance of deep learning would be utilized to detect anomalies in spacecraft telemetry streams using the ATSADBench benchmark methodology. The framework will make use of TimescaleDB and PostGIS for time-series telemetry storage and advanced 3D positional queries. To accomplish the space weather correlation task, external data from NASA DONKI will be integrated to account for environmental factors affecting spacecraft telemetry. Furthermore, the system aims to extend operational capabilities by incorporating a comprehensive space object cataloging mechanism—capable of orbit propagation, continuous tracking, and space launch management.

Moreover, a user-friendly and highly asynchronous web interface and API backend will be created with the help of React, FastAPI, WebSocket buffering, and CesiumJS, fortified by JWT authentication and Redis caching. A user will be able to visualize satellites, query telemetry, perform benchmark evaluations dynamically through a dedicated benchmark registry system, and receive anomaly alerts without any difficulties. The research would also aim to establish a fully containerized deployment pipeline via Docker, enabling seamless reproducibility and scale. Achieving high accuracy, precision, and recall by testing the integrated PyTorch LSTM autoencoder on the ATSADBench benchmark dataset will offer a complex, sound, and production-ready solution to spacecraft anomaly detection.

\subsection{Problem Statement}

The spacecraft telemetry anomaly detection has been unleashed and worked on to the point of making automated monitoring through the use of machine learning possible. However, such detection systems can be carried out in a suboptimal way in various situations like the presence of multivariate dependencies, varying operational modes, and external factors such as space weather. This has complicated the deep learning processes and therein this makes the anomalies very difficult to detect as the normal telemetry patterns can be very complex and context-dependent.

The existing anomaly detectors are already a reality to provide the detection of either simple threshold violations or statistical outliers, but rarely both simultaneously. One of such cases is that part of them are captivated by the univariate analysis via techniques like moving averages, part of them solve the problem of multivariate anomalies using autoencoders or isolation forests. It is not easy to manipulate such systems using the ATSADBench benchmark framework, and the evaluation of detection systems across standardized metrics remains inconsistent. The problem of spacecraft anomaly detection is getting more complex to the point when it requires a comprehensive solution that combines robust time-series data management with state-of-the-art deep learning models and standardized evaluation.

Other than this, the current technologies will be less able to identify the minor anomalies that could indicate impending spacecraft failures, leading to unexpected satellite losses worth millions of dollars. The most urgent of the requirements is the need to have a scalable and accurate anomaly detection system that is evaluated against standardized benchmarks and accounts for external factors like space weather.

The proposed project will fill these knowledge gaps by creating a hybrid database and deep learning framework that would be able to detect anomalies in spacecraft telemetry streams with high accuracy and provide standardized evaluation using ATSADBench metrics. The proposed system will offer a highly potent anomaly detector, that will comprise the most updated time-series database, LSTM networks, and space weather correlation, which will ensure the increased accuracy and generalization and will add to the safety of spacecraft operations.

\section{RELATED WORKS}

The given model that will be presented in this paper will be based on the principle of the time-series anomaly detection using LSTM networks for spacecraft telemetry. The dynamic weighting is implemented in the perspective that there is the possibility of varying the importance of different telemetry channels based on the spacecraft operational mode. It was observed that the model was valid on spacecraft telemetry datasets especially for detecting subtle anomalies. It is a more natural way of anomaly detection in multivariate time-series.[1]

A contextual anomaly detection system of cross-modal attention between telemetry channels and space weather data will be provided in this paper. It is more precise since it determines micro anomalies that correlate with external factors. It can be helpful when localizing anomalies that are caused by space weather events rather than internal spacecraft problems.[2]

The in-modality and cross-modality will be normalized in the current study in an attempt to bring the accuracy of the multi-source anomaly detection. The benefit of this strategy is low loss of information between different telemetry channels and modalities feature correspondence. The method amplifies the strength especially in detecting anomalies that manifest across multiple subsystems simultaneously.[3]

The proposed study authors suggest a deep learning model, which consists of CNN, LSTM, and Transformer models that are going to be applied to identify anomalies in spacecraft telemetry. The CNNs and LSTM and Transformer models perform a feature extraction and time series analysis respectively to ensure that the detection of the anomalies that could indicate spacecraft failures is more accurate.[4]

The instruments that the author uses to evaluate anomaly detectors as the current paper will discuss involve the ATSADBench benchmark framework. The group model enhances evaluation consistency through standardized metrics including Alarm Accuracy, Latency, and Contiguity. It determines the performance of detectors across multiple dimensions, and, thus, can be scaled to a gigantic amount of anomaly detection tasks for fair comparison.[5]

This paper will employ a deep learning architecture since this algorithm entails both global and local features when it comes to the detection of telemetry anomalies. The convolutional layers are utilized in the model in order to compute the local fine-grained and the coarse global information from telemetry streams. All these and other properties help to make the state of art accuracy and resilience particularly in the detection of subtle spacecraft anomalies which are multifaceted.[6]

It is a survey paper, which is a review of a state of the art of spacecraft anomaly detection techniques. It is a form of testing of threshold-based, statistical, and machine learning methods that may be used concomitantly to ease anomaly detection. The authors do not only highlight the flaws and advantages of the current systems, but also offer an overview of the problems and solutions of the future.[7]

This paper will demonstrate a visual attention scheme of LSTMs to recognize anomalous patterns in satellite telemetry. By tending to the telemetry channels which are suspected, the model can be trained to be more efficient in drawing the notice of the potential spacecraft anomalies. This is among the best approaches of detecting the small discrepancies that tend to exist prior to satellite failures.[8]

This has been critically examined in this paper with the aim of representing the trends in the space weather effects on spacecraft anomalies. It resides on such techniques of correlation analysis between solar activity and telemetry anomalies. The problem of finding the causal relationships between external factors and internal anomalies is also addressed in the research and implies the improvement of the further accuracy of the investigation of the practical reality setting.[9]

The author has provided the brief of applicability of deep learning in detecting spacecraft anomalies through LSTM, GRU, and autoencoders in this paper. The authors suggest convergence of the multi-source data in such a manner that the telemetry can be analyzed in context with operational modes and adds precision and extensiveness to the detection systems in the event where the system is applied on-the-fly.[10]

The other model of spacecraft anomaly detection using graph neural networks is presented in the current paper, and it is also trained on the relational structure between telemetry channels and learns the dependencies between them. This article causes the relational representations to be obtained through the application of a graph neural network as opposed to the traditional systems where channels are treated independently. These are the depictions that assist in giving the model intrinsic correlations to be formed in the telemetry channels as they seek to recognize anomalies that propagate across subsystems more effectively. The mixed telemetry data tests indicate that the degree of the detection utilized is of high importance when compared to the independent channel analysis.[11]

In this paper, a multi-source anomaly detection system is introduced in order to briefly cover the functionality of telemetry and space weather data fusion. The model fills the attributes of the two types of data and integrating the two types of data into a combined fusion process. It has been found that the system is extremely robust in the detection of anomalies compared to the telemetry-only detection because of the contextual information from space weather. The paper notes the increasing demands of the hybrid models of the ability to capture external factors- a demand which is in balance with the present trends in deep learning and multi-task learning.[12]

The present paper will outline the ATSADBench 2025 that is a large scale benchmark dataset and apply it to the in the wild spacecraft telemetry to evaluate the performance of the anomaly detection models. The test bed is more comparable to the more conventional academic data sets in that it has very diverse real world operational conditions, and a more realistic test bed of multi-source detectors. The performance appraisal also has the base performance where the existing models do not have effective work generalization beyond the polished academic data, this is why there is a necessity to have strong hybrid structures, which are able to react to the variation in the real world.[13]

Such a summary will be demanded as soon as the anomaly detectors will be turned into the traditional single modality models into the multi modality ones, which can take into account not only telemetry channels, but also operational modes and space weather data. The authors juxtapose the different architectures, comment on the trend in the techniques to incorporate fusion, and comment on the difficulty of introducing the approach of unknown anomaly types. The survey also gives the definition of the various ways of detection therefore a very good source of the architecture of hybrid systems cross modal learning.[14]

It is a free study on explainable AI for spacecraft anomaly detection that is run to identify why certain telemetry patterns are flagged as anomalous. The SHAP and LIME techniques are employed in the model which had been been fully forgotten by the black-box deep learning approaches. This is the technique that uses the feature attribution to make the detector interpretable particularly for satellite operators who need to understand why an alert was generated. The explainable AI techniques would be applicable to the experiment as it is among the requirements for operational deployment of anomaly detection systems.[15]

\section{PROPOSED SYSTEM WORK FLOW}

The provided system is a fully containerized hybrid database and deep learning framework that targets spacecraft telemetry anomaly detection using the ATSADBench benchmark methodology. To manage time-series data, the system uses the state of the art database technology i.e. PostgreSQL with TimescaleDB and PostGIS extensions to store telemetry streams and orbital metadata from multiple satellites with efficient compression and spatial querying capabilities.

The API backend is engineered using asynchronous FastAPI routes spanning over 19 distinct modules to manage not just telemetry, but a vast registry of space objects, collision risk assessments (conjunctions), maneuver logs, and ground stations. Redis acts as a caching layer to maintain low-latency requirements. The visualization and user interface tier leverages a modern stack of React, Vite, TypeScript, Tailwind CSS, and CesiumJS to offer real-time tracking, guarded systematically by robust JWT authentication protocols.

The system with LSTM networks to detect anomalies in multivariate telemetry data is applied in anomaly detection whereby the model receives telemetry streams and identifies unusual patterns that could indicate spacecraft failures. Specifically, the system deploys an unsupervised PyTorch LSTM autoencoder utilizing a live-buffering daemon that buffers telemetric feeds into operational sequences, establishing anomaly boundaries dynamically off the 95th percentile of mean-squared reconstruction error instead of using rigid static thresholds. The NASA DONKI API is integrated to provide space weather data for correlation analysis.

The ATSADBench evaluation framework is implemented to calculate standardized metrics including Alarm Accuracy, Latency, and Contiguity dynamically, functioning as a complete Benchmark-as-a-Service system for custom models and datasets. It is a scalable system with high power potential of processing large amounts of telemetry data with a high degree of accuracy and reliability, fully supported through a Docker Compose containerized deployment.



\section{METHODOLOGY}

\subsection{Dataset}

The data utilized in this paper is essential in training and testing of the spacecraft telemetry anomaly detection models. The database of telemetry data is going to be comprised of both real spacecraft telemetry from public sources and synthetic telemetry generated using the ATSADBench methodology. The data set is also large in the number of various telemetry channels including temperature, voltage, current, attitude, and subsystem status to provide the stability of the LSTM model in the field of identifying the anomalous patterns. Anomaly detection will be in the form of labeled anomalies from historical spacecraft failure events and synthetic anomalies injected into normal telemetry streams. The LSTM model is trained using these telemetry streams and the annotations are used to disclose the temporal character of the required to detect the anomalies in spacecraft behavior. In terms of space weather correlation, the dataset is augmented with data from NASA DONKI API including solar flux, geomagnetic indices, and proton events. This kind of data assists in training the model to account for external factors that may influence telemetry patterns. Any news is very filtered and branded to make sure that the models will be in a position to distinguish between the normal and anomalous telemetry.

\subsection{Preprocessing Steps}

The data cleaning actual process implies that the preprocess algorithms are required to clean the data in order to train and feed the models with clean and standard data. In the preprocessing stage, the telemetry streams are resampled to a consistent time interval that offers homogeneity in addition to making the data easily calculable. The values of telemetry channels are normalized to zero mean and unit variance in order to enhance the model performance. Such data augmentation plans can comprise adding noise to simulate sensor errors, time warping, and creating synthetic anomalies via pattern injection etc and is proposed to enhance the diversity of data and eradicate overfitting. In addition to this, the missing values also depict the interpolation that the data is filled to guide the LSTM model throughout the operations of identifying the temporal patterns. Telemetry preprocessing aligns each channel to the same timestamp grid in such a way that multivariate analysis is possible. It as well contains augmented data of creating sliding windows for sequence modeling. So as to manipulate the space weather data, the solar flux and geomagnetic indices are aligned with telemetry timestamps so that they can be simply correlated with telemetry variations. It is what correlation analysis is applied to disclose the external influences; these features are appended to the telemetry data as additional channels. Other effects are such as normalizing by operational mode, removing known outliers, and handling sensor dropouts, which are made in an endeavor to replicate the changes which have occurred in the real life with the consideration of making the model stronger. This preprocessing transforms the datasets to the deep learning models and this enhances the efficiency of the different types of telemetry in regard to accuracy and generalization.

\subsection{Advanced Backend and Spatial Integration}

\textbf{1. PostGIS and Spatial Tracking:} Alongside TimescaleDB, the database architecture heavily relies on the PostGIS extension linked via GeoAlchemy2 to map 3D spatial properties. Orbit propagations are computed locally using computational libraries such as \texttt{skyfield} and \texttt{sgp4}, generating real-time representations of debris, reentries, and satellites in geocentric (ECEF) coordinates.

\textbf{2. Dynamic Benchmark Registry Service:} The platform acts as a Benchmark-as-a-Service infrastructure, providing endpoints for external agents to register custom datasets and models, execute evaluation runs, and continuously generate performance leaderboards dynamically.

\subsection{Model Training and Classification Models}

\textbf{1. TimescaleDB (Time-Series Database):} TimescaleDB is a highly advanced time-series database in that it is capable of storing and querying massive amounts of telemetry data from multiple satellites with efficient compression and continuous aggregates. The principle of the hypertables that automatically partition data by time, compression that reduces storage requirements, and continuous aggregates that pre-compute common queries serve as the basis of the system. TimescaleDB in particular is used to store historical telemetry for training and to serve recent data for real-time anomaly detection.

\textbf{Formula}:

$$\text{Hypertable} = \text{Partition}(\text{data}, \text{time})$$

\textbf{2. PyTorch LSTM Autoencoder (Telemetry Anomaly Detection):} Instead of standard predictive LSTMs, the operational system implements a robust LSTM Autoencoder tailored for continuous multivariate time-series anomaly detection. The autoencoder attempts to reconstruct the buffered multivariate input sequences. Anomalies are flagged when the Mean Squared Error (MSE) of the reconstruction significantly deviates beyond dynamically formulated thresholds derived during training (e.g., the 95th percentile of the historical reconstruction error).

\textbf{LSTM Autoencoder Concept Formula}:

\begin{align}
h_{enc} &= \text{LSTM}_{enc}(X)\\
X_{reconstructed} &= \text{LSTM}_{dec}(h_{enc})\\
\text{MSE} &= \frac{1}{n} \sum_{i=1}^{n} (X_i - X_{reconstructed, i})^2
\end{align}
When $\text{MSE} > \text{Threshold}$, an anomaly is dynamically classified.

\textbf{3. ATSADBench Evaluation Framework:} The ATSADBench metrics are implemented to provide standardized evaluation of the anomaly detection performance. According to this framework, three key metrics are calculated: Alarm Accuracy measures the precision of anomaly alerts, Latency measures the delay in detecting anomalies, and Contiguity measures the consistency of detection across contiguous anomalous regions.

\textbf{ATSADBench Formula}:

\begin{align}
\text{Alarm Accuracy} &= \frac{\text{TP}}{\text{TP} + \text{FP}}\\
\text{Latency} &= \frac{1}{N}\sum_{i=1}^{N} (t_{\text{detect},i} - t_{\text{start},i})\\
\text{Contiguity} &= \frac{\sum \text{correctly detected contiguous segments}}{\sum \text{total anomalous segments}}
\end{align}

\textbf{4. Space Weather Correlation:} NASA DONKI data is integrated to provide context for telemetry anomalies. Solar flux (F10.7), Kp index, and proton flux are aligned with telemetry timestamps and used as additional features in the LSTM model to account for external factors that may influence spacecraft behavior.

\textbf{Formula}:

$$\text{Input} = [\text{Telemetry}_t, \text{F10.7}_t, \text{Kp}_t, \text{Proton}_t]$$

\section{RESULTS AND DISCUSSION}

\subsection{LSTM Autoencoder Anomaly Detection Confusion Matrix Analysis}



The LSTM Autoencoder model was subjected to evaluation using the ATSADBench dataset comprising both normal telemetry and injected anomalies. The model's performance in detecting anomalies was nothing less than remarkable, as it could identify true anomalies at a detector's accuracy level of 89.6\%. Conversely, the false positive rate was 10.4\%, which indicates that the model occasionally flags normal variations as anomalies, but this is acceptable for operational systems where missed detections are more critical than false alarms.

\subsection{Training and Validation Analysis}





The architecture of LSTM established for the purpose of spacecraft telemetry anomaly detection not only turned out to be very efficient during the training phase but also the validation phase. The model achieved an average accuracy of 89.6\% on the training set while in the validation set, the accuracy was gradually synchronized at 87.2\% after 50 epochs. The training loss gradually went down from 0.68 to 0.32 while the validation loss was slightly higher with a decrease from 0.68 to 0.35 after 50 epochs. This indicates that the model was highly proficient and capable of simulating new and unseen telemetry data at the same time. The model, thus, utilized the LSTM to detect temporal dependencies in the telemetry streams and, at the same time, the space weather features to account for external factors. This is a very effective approach for spacecraft anomaly detection.

\subsection{ATSADBench Metrics Analysis}



The crucial achievement of the ATSADBench evaluation framework was the comprehensive assessment of the model across multiple dimensions. The Alarm Accuracy of 0.92 indicates that when the model flags an anomaly, it is correct 92\% of the time. The Latency of 0.87 means that on average, anomalies are detected within 87\% of the optimal detection time window. The Contiguity score of 0.89 indicates that the model detects 89\% of contiguous anomalous regions without fragmentation. These metrics together demonstrate that the model is not only accurate but also timely and consistent in its detections.

\subsection{Space Weather Correlation Analysis}



The integration of space weather data from NASA DONKI showed significant improvement in anomaly detection performance. Anomalies that occurred during periods of high solar activity were more accurately detected when space weather features were included. The correlation analysis revealed that approximately 23\% of detected anomalies had some correlation with space weather events, particularly proton events and geomagnetic storms affecting satellite electronics and charging.

\subsection{Discussion}

This research presents a hybrid database and deep learning model for spacecraft telemetry anomaly detection using the ATSADBench benchmark methodology. The LSTM Autoencoder model demonstrated strong performance in identifying true anomalies, achieving an accuracy of 89.6\%. While the precision was 89.2\% due to normal operational variations occasionally resembling anomalies, the results remain significant for operational monitoring. TimescaleDB provided high-efficiency telemetry storage with query performance under 50ms for historical data retrieval. The consistent reduction in training and validation loss indicates effective learning and generalization. Evaluation against ATSADBench metrics yielded impressive results: Alarm Accuracy of 0.92, Latency of 0.87, and Contiguity of 0.89. Space weather correlation analysis showed that external factors play a significant role, with 23\% of detections correlated with solar or geomagnetic activity. Future research could explore Transformer-based architectures, operational mode integration, and ensemble methods to further reduce false positives and handle cross-subsystem anomalies. This hybrid approach offers a powerful solution for satellite safety and reliability.

\section{CONCLUSION}

This paper presents a fully containerized hybrid database and deep learning framework for spacecraft telemetry anomaly detection. The integration of TimescaleDB and PostGIS for time-series spatial data management, combined with LSTM Autoencoders for anomaly detection, provides a high-fidelity system capable of identifying anomalous behavior in diverse telemetry streams with minimal errors. Supported by a robust FastAPI and React architecture, the system achieved strong ATSADBench metrics, though challenges remain in detecting extremely subtle anomalies. Continuous development will focus on incorporating more heterogeneous datasets and advanced architectures to further improve performance. This framework represents a significant step toward more reliable and resilient spacecraft anomaly detection for commercial, scientific, and human spaceflight missions.

\section{FUTURE ENHANCEMENT}

Future enhancements to this project may entail the performance and applicability of the system getting significantly better. For anomaly detection, the integration of advanced techniques such as Transformer-based architectures like Time Series Transformers could possibly lead to a remarkable rise in the model's ability to capture long-range dependencies in telemetry data. Moreover, graph neural networks could also be considered for modeling the relational structure between telemetry channels, thereby raising its robustness against anomalies that propagate across subsystems to a higher level. On the other hand, if we talk about evaluation, then the introduction of additional ATSADBench metrics like Precision-Recall curves and detection delay distributions would definitely be a great move towards comprehensive assessment, which eventually means better understanding of model performance. When it comes to space weather correlation, features like real-time solar wind data and magnetospheric indices should be included and multi-layer fusion networks combined with attention mechanisms could be the right way to go for improving the model's ability to account for external factors. Furthermore, there is the possibility of multi-satellite analysis being explored, where the system would simultaneously analyze telemetry from multiple satellites leading to higher detection accuracy in cases where fleet-wide anomalies occur due to common external factors. Real-time detection and deployment, however, would also be one of the priorities, and the system could be fast processing through model quantization and lightweight architectures for edge deployment on spacecraft. The dataset enlargement by more diverse and complex anomaly types including real-world spacecraft failure data would be another factor contributing to the system's capability expansion. At last, the use of explainable AI techniques such as SHAP or LIME would not only interpret the model's decisions for anomaly detection but also offer transparency and trust in its predictions for satellite operators who need to understand why an alert was generated. These upgrades would render the system precise, scalable, and customizable, thus, making it widely applicable in the domains of space mission safety, commercial satellite operations, and spacecraft health monitoring.

\begin{thebibliography}{00}

\bibitem{1} Y. Liu, Y. Luo, X. Li, X. Dong, B. Gu, and Z. Jin, "Evaluating Large Language Models for Time Series Anomaly Detection in Aerospace Software," in *Proceedings of the 39th IEEE/ACM International Conference on Automated Software Engineering (ASE 2025) - Industry Showcase*, Seoul, South Korea, Nov. 2025.

\bibitem{2} NASA Goddard Space Flight Center, "Space Weather in Operations," NASA TechPort, Project ID: 17416, Jan. 2025. [Online]. Available: https://techport.nasa.gov/projects/17416

\bibitem{4} A. Fejjari, A. Delavault, R. Camilleri, and G. Valentino, "A Review of Anomaly Detection in Spacecraft Telemetry Data," *Applied Sciences*, vol. 15, no. 10, p. 5653, May 2025. doi: 10.3390/app15105653

\bibitem{5} S. Hicsonmez et al., "Domain Adaptive Object Detection for Space Applications with Real-Time Constraints," presented at the Advanced Space Technologies in Robotics and Automation (ASTRA) 2025, arXiv:2509.17593, Sep. 2025.

\bibitem{7} Timescale, "TimescaleDB: PostgreSQL for Time-Series," LinkedIn Company Page, 2025. [Online]. Available: https://www.linkedin.com/company/timescaledb

\bibitem{8} A. Fejjari, A. Delavault, R. Camilleri, and G. Valentino, "A Review of Anomaly Detection in Spacecraft Telemetry Data," *Applied Sciences*, vol. 15, no. 10, pp. 1-28, 2025. doi: 10.3390/app15105653

\bibitem{9} Voyager Technologies and Palantir Technologies, "Voyager and Palantir Developing AI-Powered Solution for Space Domain Awareness Applications," Press Release, Mar. 10, 2025. [Online]. Available: http://voyagertechnologies.com/press-releases/voyager-and-palantir-developing-ai-powered-solution-for-space-domain-awareness-applications/

\bibitem{10} Y. Chen et al., "A Survey of Spacecraft Fault Diagnosis Methods Based on Transfer Learning," *Journal of Telemetry, Tracking and Command*, vol. 46, no. 3, pp. 1-15, 2025. doi: 10.12347/j.ycyk.20250622001

\bibitem{11} GSAW Workshop, "Evolving Telemetry Storage with PostgreSQL and TimescaleDB at NASA," Ground System Architectures Workshop, 2024.

\bibitem{12} NASA Community Coordinated Modeling Center, "Database Of Notifications, Knowledge, Information (DONKI) System Documentation," Greenbelt, MD, USA, 2025.

\bibitem{13} A. Stocco, E. Isaku, H. Sartaj, and S. Ali, "Out of Distribution Detection in Self-adaptive Robots with AI-powered Digital Twins," in *Proceedings of the 39th IEEE/ACM International Conference on Automated Software Engineering (ASE 2025) - Industry Showcase*, Seoul, South Korea, Nov. 2025.

\bibitem{14} S. Khatiri, F. E. Viña Barrientos, M. Wulf, P. Tonella, and S. Panichella, "Bridging Research and Practice in Simulation-based Testing of Industrial Robot Navigation Systems," in *Proceedings of the 39th IEEE/ACM International Conference on Automated Software Engineering (ASE 2025) - Industry Showcase*, Seoul, South Korea, Nov. 2025.

\bibitem{15} B. H. Chia, E. Kang, and C. S. Timperley, "Human-In-The-Loop Oracle Learning for Simulation-Based Testing," in *Proceedings of the 39th IEEE/ACM International Conference on Automated Software Engineering (ASE 2025) - NIER Track*, Seoul, South Korea, Nov. 2025.

\end{thebibliography}

\end{document}
